\section{Signed source response and displacement of a double event}
\label{s:response}

The two-root flow separates the analytic geometry of a possible collision
from the source information needed to exclude it.  This section gives a
signed source functional, an explicit nonlinear error, and a continuation
criterion for the event itself.  The results consolidate the analytic
arguments in \cite{note08}; their proofs are included here.

\subsection{An exact response formula at a fixed real time}

Fix a real time, which is suppressed in the notation, and a real spatial
center $a$.  Let $F$ and $H=F+E$ be holomorphic near
$\overline{B(a,R)}$.  Assume that $F$ is nonzero on
$\Gamma=\{z:|z-a|=R\}$, has exactly two zeros counted with multiplicity
inside $\Gamma$, and that both zeros belong to $B(a,r)$, where $0<r<R$.
Write $c$ for their centroid and $\Delta_F$ for their squared difference.
The reference zeros may coincide.  Throughout this section contours have
counterclockwise orientation, and
\[
 \mathcal I[g]=\frac{1}{2\pi i}\int_\Gamma g(z)\,dz.
\]
Suppose the complete wall estimate is
\begin{equation}
 \rho:=\sup_\Gamma|E/F|<1.
 \label{s:ratio}
\end{equation}
Rouch\'e's theorem then gives two zeros of $F+sE$ inside $\Gamma$ for
$0\le s\le1$.  A count in a smaller disk, if needed for the two-root
flow theorem, is an additional condition; it does not follow from
\eqref{s:ratio} alone.

\begin{theorem}[Logarithmic discriminant response]
\label{s:log-response}
Let $\ell=\log(1+E/F)$ be the branch defined by its convergent power
series on an annulus containing $\Gamma$.  Then
\begin{equation}
 \Delta_H-\Delta_F
 =-4\mathcal I[(z-c)\ell]-\bigl(\mathcal I[\ell]\bigr)^2.
 \label{s:exact-response}
\end{equation}
Set
\begin{align}
 \mathcal L_F(E)&=-4\mathcal I[(z-c)E/F],\label{s:linear-response}\\
 l_1(\rho)&=-\log(1-\rho),\qquad
 l_2(\rho)=-\log(1-\rho)-\rho,\notag\\
 \mathcal B(\rho)&=4R(R+r)l_2(\rho)+R^2l_1(\rho)^2.
 \label{s:quadratic-payment}
\end{align}
The finite response satisfies
\begin{equation}
 |\Delta_H-\Delta_F-\mathcal L_F(E)|\le\mathcal B(\rho).
 \label{s:response-bound}
\end{equation}
In particular the error is $O(\rho^2)$, with constants independent of
the distance between the two reference zeros.
\end{theorem}

\begin{proof}
Put $w=z-a$, and for $G=F,H$ let
$S_k(G)=\mathcal I[w^kG_z/G]$ be the sum of the $k$th powers of its
two centered zeros, counting multiplicity.  On the annulus,
$H_z/H-F_z/F=\ell_z$.  Integration by parts on the closed contour gives
\[
 S_1(H)-S_1(F)=-\mathcal I[\ell],\qquad
 S_2(H)-S_2(F)=-2\mathcal I[w\ell].
\]
The identity $\Delta_G=2S_2(G)-S_1(G)^2$, together with
$S_1(F)=2(c-a)$, proves \eqref{s:exact-response}.  No extension of
$\ell$ across the zeros of $F$ was used.  Its power series yields
$|\ell|\le l_1(\rho)$ and $|\ell-E/F|\le l_2(\rho)$ on $\Gamma$.
Now $|\mathcal I[g]|\le R\sup_\Gamma|g|$ and
$|z-c|\le R+r$ prove \eqref{s:response-bound}.
\end{proof}

If $F$ and $H$ are real-symmetric, their discriminants and
$\mathcal L_F(E)$ are real.  Thus a signed enclosure
$\mathcal L_F(E)\in[L_-,L_+]$ gives
\begin{equation}
 \Delta_F+L_- -\mathcal B(\rho)
 \le\Delta_H\le
 \Delta_F+L_+ +\mathcal B(\rho).
 \label{s:signed-enclosure}
\end{equation}
A positive lower bound at a negative real cap can be combined with the
two-root charge theorem once all its spatial and phase hypotheses have
been established.  A positive lower bound at $t=0$ itself excludes an
endpoint double in this fiber.  Formula \eqref{s:signed-enclosure}
does not provide either strict inequality without further source
information.  It also retains the external wall ratio and the radii;
being finite at an internal double is not a uniform condition estimate
over all heights.

\subsection{The complete signed source functional}

Suppose, now displaying the fixed real time $t$, that
\[
 E(t,z)=\int_0^\infty\psi(u)e^{tu^2}\cos(zu)\,du.
\]
If $|F(t,z)|\ge\delta>0$ on $\Gamma$ and
\begin{equation}
 N_\psi(t,R):=\int_0^\infty|\psi(u)|e^{tu^2}\cosh(Ru)\,du<\infty,
 \label{s:complete-residual-norm}
\end{equation}
then Fubini's theorem gives
\begin{align}
 \mathcal L_F(E)&=\int_0^\infty\psi(u)e^{tu^2}K_F(u)\,du,
 \label{s:source-functional}\\
 K_F(u)&=-4\mathcal I\!\left[
             \frac{(z-c)\cos(zu)}{F(t,z)}\right],
 \qquad
 |K_F(u)|\le\frac{4R(R+r)}{\delta}\cosh(Ru).
 \label{s:source-kernel}
\end{align}
Indeed $|\operatorname{Im}z|\le R$ on the circle because $a$ is real,
so $|\cos(zu)|\le\cosh(Ru)$.  This proves an integrable majorant for
the exchanged integrals, and also
$\rho\le N_\psi(t,R)/\delta$.  Conjugation symmetry makes $K_F$ real
on the positive real $u$ axis.

At an exact double $b=c$, factor $F(t,z)=U(t,z)(z-b)^2$ with $U$ a
nonvanishing analytic function.  A residue in
\eqref{s:linear-response} then gives the particularly simple identity
\begin{equation}
 \mathcal L_F(E)=-\frac{4E(t,b)}{U(t,b)}
               =-\frac{8E(t,b)}{F_{zz}(t,b)}.
 \label{s:double-response}
\end{equation}
The sign is therefore an oscillatory source measurement divided by a
nonzero curvature.  Positivity of the source and absolute moment bounds
alone do not determine it.  For numerical use, every omitted source
mode, unbounded integration ray, contour arc, and quadrature error must
be included in both the wall ratio and the signed functional.  No new
numerical acquisition is used to prove this identity.

\subsection{Differentiating a real event through an affine heat family}

Let $H_s(t,z)=F(t,z)+sE(t,z)$ be real-symmetric, with both $F$ and $E$
satisfying the backward heat equation.  Suppose $F$ has an ordinary real
double at $(\tau_0,b_0)$, so $F=F_z=0$ and $F_{zz}\ne0$ there.

\begin{proposition}[Event derivatives and a one-sided finite error]
\label{s:event-derivatives}
There is a unique local real-analytic event branch
$s\mapsto(\tau_s,b_s)$.  At a point on this branch, set
\[
 N=H_{s,zz}(\tau_s,b_s),\quad M=H_{s,zzz}(\tau_s,b_s),\quad
 P_4=H_{s,zzzz}(\tau_s,b_s).
\]
With the residual derivatives evaluated at the same point,
\begin{align}
 \tau_s'&=\frac{E}{N}, &
 b_s'&=\frac{M\tau_s'-E_z}{N},\label{s:event-first}\\
 \tau_s''&=\frac{P_4}{N}(\tau_s')^2
          -\frac{2E_{zz}}{N}\tau_s'-(b_s')^2.
 \label{s:event-second}
\end{align}
If the branch exists for $0\le s\le1$ and along it
\[
 |N|\ge\nu>0,\quad |P_4|\le C_4,\quad |E|\le E_0,
 \quad |E_{zz}|\le E_2,
\]
then
\begin{equation}
 \tau_1\le\tau_0+\frac{E(\tau_0,b_0)}{F_{zz}(\tau_0,b_0)}
       +\frac12\left(\frac{C_4E_0^2}{\nu^3}
                         +\frac{2E_2E_0}{\nu^2}\right).
 \label{s:event-location}
\end{equation}
\end{proposition}

\begin{proof}
The Jacobian of $(H_s,H_{s,z})$ with respect to $(t,z)$ at an event is
\[
 \begin{pmatrix}-N&0\\-M&N\end{pmatrix},
\]
whose determinant is $-N^2$.  The analytic implicit-function theorem
gives the local branch, and real symmetry gives its reality.  One
derivative of the two event equations gives
$-N\tau_s'+E=0$ and $-M\tau_s'+Nb_s'+E_z=0$.
For a second derivative of the first equation, use
$H_{s,tt}=P_4$, $H_{s,tz}=-M$, and $E_t=-E_{zz}$ to obtain
\[
 0=-N\tau_s''+P_4(\tau_s')^2-2M\tau_s'b_s'
       +N(b_s')^2-2E_{zz}\tau_s'+2E_zb_s'.
\]
Substitution of $E_z=M\tau_s'-Nb_s'$ leaves
\eqref{s:event-second}.  All terms are real.  In particular
$-(b_s')^2\le0$, whereas $|\tau_s'|\le E_0/\nu$, so
\[
 \tau_s''\le\frac{C_4E_0^2}{\nu^3}
                         +\frac{2E_2E_0}{\nu^2}.
\]
Integrating twice on $[0,1]$ proves \eqref{s:event-location}.
\end{proof}

This calculation also checks the signs in the discriminant formula.
At a double, $\partial_s\Delta_s=-8E/N$ by
\eqref{s:double-response}, while $\partial_t\Delta_s=8$ by the
two-root heat equation.  Differentiating
$\Delta_s(\tau_s)=0$ recovers $\tau_s'=E/N$.

\subsection{Paying continuation and curvature from common walls}

The local implicit-function theorem by itself does not continue an
event to $s=1$.  The following conditions supply that missing step and
the curvature floor in \eqref{s:event-location}.

\begin{proposition}[A complete conditional continuation criterion]
\label{s:continuation}
Let $\mathcal D$ be a closed time rectangle with real caps $t_-<t_+$
and imaginary half-width $\tau$.  Suppose every $H_s=F+sE$,
$0\le s\le1$, has exactly two zeros inside $B(a,r)$ throughout
$\mathcal D$, none in the closed collar to $\Gamma$, a common outer
wall floor $|H_s|\ge\delta>0$, and a common strict phase condition
from the two-root real-cap theorem.  Assume
\[
 \Delta_F(t_-)<0<\Delta_F(t_+),\qquad
 m_{\rm cap}=\min\{-\Delta_F(t_-),\Delta_F(t_+)\}>0.
\]
If the two outer cap circles satisfy $|E/F|\le\rho<1$ and
\begin{equation}
 C_{\rm cap}:=4R(R+r)l_1(\rho)+R^2l_1(\rho)^2<m_{\rm cap},
 \label{s:cap-persistence}
\end{equation}
then a unique real event exists for every $s\in[0,1]$, the events form
one real-analytic branch, and
\begin{equation}
 |H_{s,zz}(\tau_s,b_s)|\ge\frac{2\delta}{(R+r)^2}.
 \label{s:curvature-payment}
\end{equation}
\end{proposition}

\begin{proof}
Apply \eqref{s:exact-response} to $sE$ at each cap.  The bound
$|\ell|\le l_1(\rho)$ gives
$|\Delta_s(t_\pm)-\Delta_F(t_\pm)|\le C_{\rm cap}$.
Thus both cap signs persist.  The two-root real-cap theorem gives
exactly one event for each $s$, located on the real time interval.  It
has spatial multiplicity two and $\partial_t\Delta_s=8$.  The local
branches of Proposition~\ref{s:event-derivatives} agree wherever their
domains overlap by uniqueness.  They extend along the full compact
parameter interval: any limit of event points remains away from the
time caps by the strict cap margins and from the spatial boundary by
the common walls, and is again a double.  The implicit-function theorem
then extends the branch at that limit.

For the curvature, write $H_s=U_sP_s$ in the disk.  On $\Gamma$,
$|P_s|\le(R+r)^2$, hence $|U_s|\ge\delta/(R+r)^2$.  Since $U_s$
has no zeros, the maximum modulus principle applied to $1/U_s$
propagates this bound throughout the disk.  At an event,
$P_s=(z-b_s)^2$ and $H_{s,zz}=2U_s$, proving
\eqref{s:curvature-payment}.
\end{proof}

Here are sufficient reference data for all the common hypotheses.
Suppose $F$ has two zeros inside each of the disks of radii $r,R$
throughout $\mathcal D$, with wall floors $\delta_{F,r},\delta_{F,R}$,
and
\[
 Q_v=\sup_{t\in\mathcal D,\ |z-a|=v}|E/F|<1\quad(v=r,R).
\]
Rouch\'e's theorem supplies both counts for every $H_s$, with floors
$\delta_v=(1-Q_v)\delta_{F,v}$.  Take $\delta=\delta_R$.  If a
complete outer-wall derivative bound is $|F_z|+|E_z|\le L$, the
sufficient common phase inequality is
\begin{equation}
 \frac{4RL}{\delta(R-r)^2}\,\tau<\frac\pi2.
 \label{s:common-phase}
\end{equation}
A more precise common estimate for the regular coefficient $A_s$ may
replace this one.  If the source densities are $\Phi_0$ and $\psi$,
the remaining bounds along the real event branch can be taken to be
\begin{align}
 E_j&=\int_0^\infty|\psi(u)|u^je^{t_+u^2}\,du\quad(j=0,2),\notag\\
 C_4&=\int_0^\infty\bigl(|\Phi_0(u)|+|\psi(u)|\bigr)
                           u^4e^{t_+u^2}\,du.
 \label{s:moment-payments}
\end{align}
The real event positions allow $|\cos(b_su)|\le1$ in these bounds.
All source modes and integration rays are included.

Consequently \eqref{s:event-location}, together with
\eqref{s:cap-persistence}--\eqref{s:moment-payments}, is a complete
conditional criterion for locating the unique event below zero.  It
does not supply a negative right-hand side for the unchanged theta
source.  Both the signed initial velocity and the external wall
margins remain substantive inputs.

\begin{example}[An exact nonlinear and continuation check]
\label{s:quadratic-control}
For
\[
 H_s(t,z)=z^2-2t+s(bz+c),
\]
one has
\[
 \Delta_s(t)=8t+s^2b^2-4sc,\qquad
 \tau_s=\frac{sc}{2}-\frac{s^2b^2}{8},\qquad b_s=-\frac{sb}{2}.
\]
The initial signed response is $-4c$, while the finite correction is
$s^2b^2$.  The logarithmic integral alone in
\eqref{s:exact-response} contributes $2s^2b^2$; its last square is
essential.  The event equations give $\tau_0'=c/2$ and
$\tau_s''=-b^2/4$, as required.

The continuation inequalities can also be checked with strict rational
margins.  Set $r=1$, $R=2$, $|\operatorname{Re}t|,
|\operatorname{Im}t|\le1/4$, $b=1/16$, and $c=\pm1/32$.
Since $2|t|\le\sqrt2/2<3/4$, the inner and outer reference wall
floors exceed $1/4$ and $13/4$.  Thus the relative errors are less
than $3/8$ and $5/104$, respectively, and both counts persist.  The
outer floor of $H_s$ exceeds $99/32$, so
\eqref{s:curvature-payment} permits $\nu=11/16$.
Here $A_s=0$ exactly, and the reference cap discriminants are $-2,+2$.
Since $l_1(5/104)\le5/99$,
\[
 C_{\rm cap}\le\frac{120}{99}+\frac{100}{9801}<2.
\]
For $c=-1/32$, one has $P_4=E_{zz}=0$, so
\eqref{s:event-location} proves $\tau_1\le-1/64<0$; its exact value
is $-1/64-1/2048$.  For $c=+1/32$ the event is positive.  The same
geometric conditions accommodate both outcomes; the sign comes from
the signed displacement.
\end{example}

\section{What finite complete envelope moments cannot determine}
\label{s:moment-obstruction}

The obstruction proved here concerns precisely the envelope moments
used in absolute source estimates.  It does not apply to an arbitrary
finite list of observations, which could include the oscillatory
measurement one wishes to determine.

\begin{theorem}[Fixed envelope moments with both event directions]
\label{s:finite-moment-control}
Fix any finite collection
\begin{equation}
 k_j(u)=u^{n_j}e^{T_ju^2}\cosh(Y_ju),\qquad
 n_j\in\mathbb Z_{\ge0},\quad T_j\in\mathbb R,\quad Y_j\ge0.
 \label{s:envelope-moments}
\end{equation}
There is a strictly positive smooth even source $\Phi_0$, with
super-Gaussian decay, and a real smooth perturbation $\psi$ supported
in a compact subinterval of $(0,\infty)$, such that for all sufficiently
small positive and negative $s$:
\begin{enumerate}
 \item $\Phi_s=\Phi_0+s\psi$ remains positive, with the same source
       near zero and outside a fixed compact interval;
 \item every selected complete moment
       $\int_0^\infty\Phi_s(u)k_j(u)\,du$ is independent of $s$;
 \item the associated heat family has a real ordinary-double branch
       with $\tau_0=0$ and $\tau_0'>0$.
\end{enumerate}
In particular the double event lies on either side of zero for the
respective signs of sufficiently small $s$.
\end{theorem}

\begin{proof}
Take $c_0>0$ and
\[
 \Phi_0(u)=c_0\left(e^{-u^4}
          +\tfrac12e^{-(u-4)^4}+\tfrac12e^{-(u+4)^4}\right).
\]
This is positive, even, and rapidly decaying.  If
$G(z)=\int_0^\infty e^{-u^4}\cos(zu)\,du$, translation of the
whole-line Fourier integral gives
\[
 H_0(0,z)=c_0G(z)(1+\cos4z).
\]
At $b=\pi/4$ the last factor has an ordinary double.  To see that
$G(b)>0$ without numerical evaluation, write
$M_j=\int_0^\infty u^je^{-u^4}\,du$.  Integration of
$(ue^{-u^4})'$ gives $M_4=M_0/4$; Cauchy--Schwarz gives
$M_2\le(M_0M_4)^{1/2}=M_0/2$.  Since
$\cos x\ge1-x^2/2$,
\[
 G(b)\ge M_0-\frac{b^2}{2}M_2
       \ge M_0(1-b^2/4)>0.
\]
Consequently $H_{0,zz}(0,b)=16c_0G(b)>0$.

Choose a bounded open interval $I\Subset(0,\infty)$ and a smooth
nonnegative bump $\chi$ supported in $I$ and positive on a subinterval.
In the real inner product
$\langle f,g\rangle_\chi=\int_0^\infty\chi fg$, let $p$ be the
orthogonal projection of $f(u)=\cos(bu)$ onto the span of the $k_j$;
linearly dependent functions may first be discarded.  Set
$\psi=\chi(f-p)$.  Then
\begin{equation}
 \int_0^\infty\psi k_j=0,
 \qquad
 \int_0^\infty\psi(u)\cos(bu)\,du
    =\int_0^\infty\chi(u)(f(u)-p(u))^2\,du>0.
 \label{s:projection-control}
\end{equation}

For the strict inequality, $\cos(bu)$ cannot belong to the indicated
finite span on any interval.  Such an equality would extend to the
whole real axis by analyticity.  Expanding each hyperbolic cosine and
grouping equal exponential factors writes its proposed right-hand side
as a finite sum of polynomial multiples of $e^{Tu^2+Yu}$ with real
$T,Y$.  As $u\to+\infty$, eliminate the largest quadratic exponent
$T>0$, and within it the largest linear exponent, since a nonzero
polynomial times that factor cannot remain bounded.  Repeating removes
all growing terms, including those with $T=0,Y>0$.  All terms with
$T<0$, or with $T=0,Y<0$, decay.  The remaining term is a polynomial,
which must be constant to remain bounded.  The resulting right-hand
side has a limit, whereas $\cos(bu)$ does not.  This contradiction
proves \eqref{s:projection-control}.

The base source has a positive minimum on $\operatorname{supp}\psi$;
therefore both signs of sufficiently small $s$ preserve positivity.
The even extension of $\psi$ is smooth because its support is away
from zero, and all prescribed moments remain exactly fixed.  The
super-Gaussian tail and the compact perturbation give jointly entire
heat transforms.  Proposition~\ref{s:event-derivatives} and
\eqref{s:projection-control}, divided by the positive curvature,
give $\tau_0'>0$.  The local implicit-function theorem supplies the
branch and the asserted signs.
\end{proof}

A sufficiently small proper two-root time-space window exists around
the base double: isolate the zero at time zero, use nonvanishing on
nested spatial circles to preserve the two counts, and shrink the time
rectangle until the regular coefficient satisfies the strict phase
condition.  Its strict margins persist for all sufficiently small
$|s|$ by uniform convergence on compact sets.  Thus the two event
directions coexist with the same selected moments and the same strict
local analytic requirements.

\begin{remark}[Scope of the control]
The preserved moment vector is that of the constructed source
$\Phi_0$, not the theta source.  Compact smooth perturbations need not
preserve analyticity in the source variable, stronger shape conditions,
or additional primitive identities.  This construction asserts no
global all-real heat slice, no all-positive-time real-rootedness, no
canonical theta coefficient expansion, and no ordinary Dirichlet
series or full normalized arithmetic reflection identity.  It excludes
only a universal inference from the displayed finite envelope moments
and strict local analytic data.  A measured signed functional such as
\eqref{s:source-functional}, exact arithmetic data, and further
consequences of RH lie outside that information class.  In particular,
this preservation list must not be combined with the distinct lists
belonging to the other collision controls.
\end{remark}

\section{The canonical theta tail in the discriminant response}
\label{s:canonical-tail}

The actual theta tail has more structure than an arbitrary positive
source bounded between two envelopes.  We construct that structure
explicitly, following the source identity in \cite{tail04}, and then
apply it to the signed kernel.  The resulting interval is strictly
narrower than a specified pointwise-envelope interval, but is not
proved to exclude zero.

\subsection{An explicit positive Laplace law}

Retain the complete mode normalization
\[
 \lambda_n=\pi n^2,\qquad
 \phi_n(u)=\lambda_ne^{5u}(2\lambda_ne^{4u}-3)
                              e^{-\lambda_ne^{4u}}.
\]
Every mode is positive for $u\ge0$.  Fix an integer $m\ge1$, and set
\begin{gather}
 S_m=\sum_{n=1}^m\phi_n,\qquad R_m=\sum_{n>m}\phi_n,\notag\\
 \alpha=\lambda_m,\quad\beta=\lambda_{m+1},\quad d=\beta-\alpha,
 \quad v=e^{4u}-1,\quad D_0=2\alpha-3,\quad
 \gamma=\frac{D_0}{2\alpha}>0.\notag
\end{gather}
For $n>m$, define the nonnegative quantities
\begin{equation}
 d_n=\lambda_n-\alpha,\qquad e_n=\lambda_n-\beta,
 \qquad
 a_n=\frac{\lambda_n^2}{\alpha^2}e^{-d_n},\qquad
 b_n=\frac{3\lambda_nd_n}{\alpha^2D_0}e^{-d_n}.
 \label{s:tail-masses}
\end{equation}

\begin{proposition}[Canonical tail measure]
\label{s:laplace-law}
Let $\mathsf X_\gamma$ have exponential density
$\gamma e^{-\gamma x}$ on $[0,\infty)$, and put
\begin{align}
 E_m&=\sum_{n>m}(a_n+b_n)=\frac{R_m(0)}{\phi_m(0)},\notag\\
 \nu_m&=\frac1{E_m}\sum_{n>m}
       \left[a_n\delta_{e_n}
                +b_n\operatorname{Law}(e_n+\mathsf X_\gamma)\right],
 \label{s:tail-probability}\\
 p_m(v)&=\int_0^\infty e^{-xv}\,d\nu_m(x).
 \label{s:laplace-transform}
\end{align}
Then $\nu_m$ is a probability measure and
\begin{equation}
 R_m(u)=E_m\phi_m(u)e^{-dv}p_m(v).
 \label{s:exact-tail}
\end{equation}
If $A=a_{m+1}$ and $\eta=A/E_m$, then
\[
 0<\eta<1,\qquad p_m(0)=1,\qquad p_m(\infty)=\eta,
 \qquad (-1)^jp_m^{(j)}(v)\ge0\quad(j\ge0,v\ge0).
\]
In particular
\begin{equation}
 A\phi_m(u)e^{-dv}\le R_m(u)\le E_m\phi_m(u)e^{-dv}.
 \label{s:tail-envelope}
\end{equation}
\end{proposition}

\begin{proof}
Division of the source polynomials, retaining their lower-order terms,
gives
\begin{equation}
 \frac{\phi_n(u)}{\phi_m(u)}
   =e^{-d_nv}\left(a_n+b_n\frac{\gamma}{\gamma+v}\right).
 \label{s:mode-division}
\end{equation}
For example, with $q=1+v$ the needed algebra is
\[
 \frac{2\lambda_nq-3}{2\alpha q-3}
 =\frac{\lambda_n}{\alpha}
    +\frac{3(\lambda_n-\alpha)}{\alpha(2\alpha q-3)}.
\]
Also $\int_0^\infty e^{-xv}\gamma e^{-\gamma x}\,dx
=\gamma/(\gamma+v)$.  Factor $e^{-dv}$ out of
\eqref{s:mode-division} and sum over all $n>m$ to obtain
\eqref{s:exact-tail} and \eqref{s:tail-probability}.  Evaluating at
$v=0$ proves the expression for $E_m$ and total mass one.  The masses
have Gaussian decay in $n$, and the continuous pieces have exponential
tails, so every moment of $\nu_m$ is finite.  Differentiation under its
integral proves the alternating derivative inequalities, including at
zero.  Its only mass at zero is $a_{m+1}/E_m$; the exponential pieces
have no atoms, and $b_{m+1}>0$, so this number lies strictly between
zero and one.  Dominated convergence gives $p_m(\infty)=\eta$.
The last enclosure follows.
\end{proof}

For later use, $E_m/A$ is bounded as $m$ varies and tends to one.  Here
is a direct estimate.  Put $k=m+1$ and
\[
 \varrho_m=\left(\frac{k+1}{k}\right)^4e^{-\pi(2k+1)}<1.
\]
Successive ratios of $a_n$ for $n\ge k$ are at most $\varrho_m$,
and $b_n/a_n=3d_n/(\lambda_nD_0)\le3/D_0$.  Consequently
\begin{equation}
 1\le\frac{E_m}{A}
 \le\frac{1+3/D_0}{1-\varrho_m}\longrightarrow1.
 \label{s:tail-mass-ratio}
\end{equation}
This also supplies a complete elementary majorant for the infinite
mass sum.

\subsection{A cumulative-integral enclosure for the actual response}

At a fixed real cap take the \emph{unshifted} reference and residual
\begin{equation}
 F(t,z)=\int_0^\infty S_m(u)e^{tu^2}\cos(zu)\,du,
 \qquad E(t,z)=\int_0^\infty R_m(u)e^{tu^2}\cos(zu)\,du.
 \label{s:unshifted-head}
\end{equation}
Assume $F$ has the proper two-root geometry of Section~\ref{s:response}
and outer wall floor $\delta>0$.  Define
\begin{align}
 W(u)&=e^{tu^2}\phi_m(u)e^{-d(e^{4u}-1)},\notag\\
 z_m(u)&=\frac{p_m(e^{4u}-1)-\eta}{1-\eta},\qquad
 d\sigma_m(w)=-z_m'(w)\,dw.
 \label{s:monotone-tail}
\end{align}
The function $z_m$ decreases from one to zero.  Hence $\sigma_m$ is
a probability measure, and $z_m(u)=\sigma_m([u,\infty))$.
The source identity becomes
\begin{equation}
 e^{tu^2}R_m(u)=W(u)\{A+(E_m-A)z_m(u)\}.
 \label{s:affine-multiplier}
\end{equation}
For the real kernel in \eqref{s:source-kernel}, put
\[
 J_K=\int_0^\infty W(u)K_F(u)\,du,
 \qquad C_K(w)=\int_0^wW(u)K_F(u)\,du,
 \qquad C_K(\infty)=J_K.
\]
The bound in \eqref{s:source-kernel} and the complete source decay
make these integrals absolutely convergent.

\begin{theorem}[Canonical signed discriminant interval]
\label{s:tail-discriminant}
With the preceding notation,
\begin{equation}
 \mathcal L_F(E)=A J_K+(E_m-A)
                        \int_0^\infty C_K(w)\,d\sigma_m(w).
 \label{s:tail-response}
\end{equation}
Suppose a complete wall-ratio payment is
\begin{equation}
 \overline\rho=\frac{E_m}{\delta}
                \int_0^\infty W(u)\cosh(Ru)\,du<1.
 \label{s:tail-ratio}
\end{equation}
Then
\begin{align}
 \Delta_H&\ge\Delta_F+A J_K+(E_m-A)\min_{0\le w\le\infty}C_K(w)
                         -\mathcal B(\overline\rho),\notag\\
 \Delta_H&\le\Delta_F+A J_K+(E_m-A)\max_{0\le w\le\infty}C_K(w)
                         +\mathcal B(\overline\rho).
 \label{s:canonical-interval}
\end{align}
The extrema exist on the compactified half-line.  A sharper established
wall ratio can replace \eqref{s:tail-ratio}.
\end{theorem}

\begin{proof}
Substitute \eqref{s:affine-multiplier} into
\eqref{s:source-functional}, write
$z_m(u)=\int\mathbf1_{u\le w}\,d\sigma_m(w)$, and use Fubini.
Absolute integrability follows from $\sigma_m$ having mass one and
$\int W|K_F|<\infty$.  This proves \eqref{s:tail-response}.
Continuity of $C_K$ and its limit at infinity give a compact range,
so its probability average lies between its minimum and maximum.
The upper envelope in \eqref{s:tail-envelope} proves
$\sup_\Gamma|E/F|\le\overline\rho$.  Finally
\eqref{s:response-bound}, with the increasing function
$\mathcal B$, proves \eqref{s:canonical-interval}.
\end{proof}

This enclosure retains every theta mode and the entire source ray.
It is an analytic estimate, independent of a sampling scheme.  An
implementation must additionally enclose the acquisition errors and
the full-domain extrema of $C_K$; no positive numerical lower endpoint
is asserted here.

\subsection{Strict improvement over a matched pointwise envelope}

The following comparison keeps the same known lower and upper source
bounds, the same wall ratio, and the same nonlinear response error.
The pointwise comparison drops only the monotonicity of $z_m$.

\begin{lemma}[Oscillation of the two-root response kernel]
\label{s:kernel-oscillation}
If both reference roots lie in $\operatorname{Re}z>0$, then $K_F$
has strictly positive and strictly negative intervals arbitrarily far
along the positive $u$ axis.
\end{lemma}

\begin{proof}
Factor $F=UP$.  For distinct roots $z_1,z_2$, their centroid satisfies
$(z_i-c)/P_z(z_i)=1/2$.  The residue theorem therefore gives
\begin{equation}
 K_F(u)=-2\left(\frac{\cos(z_1u)}{U(z_1)}
                         +\frac{\cos(z_2u)}{U(z_2)}\right).
 \label{s:kernel-residues}
\end{equation}
At a double $b>0$ the formula is
$K_F(u)=-4\cos(bu)/U(b)$, which has the asserted signs.
For distinct real positive roots, \eqref{s:kernel-residues} is a
nonzero real combination of two distinct positive-frequency cosines.
Its primitive is bounded and its mean square is positive: integrating
the cosine products shows that the time average of the square tends
to half the sum of the squared coefficients.  If the kernel were
eventually nonnegative, its integral on the remaining ray would be
bounded, whereas $K_F^2\le\|K_F\|_\infty K_F$ would then bound the
square integral, contradicting its positive linear growth.  Eventual
nonpositivity is excluded in the same way.  Continuity converts the
resulting arbitrarily distant points of either sign into intervals.

The other possibility under real symmetry is a conjugate pair
$b\pm iy$ with $b,y>0$.  Write $1/U(b+iy)=c+ie$.  Equation
\eqref{s:kernel-residues} becomes
\[
 K_F(u)=-4\{c\cos(bu)\cosh(yu)+e\sin(bu)\sinh(yu)\}.
\]
After division by $e^{yu}$ this is a nonzero sinusoid plus a term
tending exponentially to zero; $(c,e)\ne(0,0)$ because $U$ is
nonvanishing.  Its positive and negative lobes give the claim.
\end{proof}

\begin{theorem}[Strict matched-envelope improvement]
\label{s:strict-improvement}
Under the hypotheses of Theorem~\ref{s:tail-discriminant}, assume also
that both reference roots lie in $\operatorname{Re}z>0$.  Set
\[
 f(u)=W(u)K_F(u),\qquad
 P_+=\int_0^\infty f_+(u)\,du,
 \qquad P_-=\int_0^\infty f_-(u)\,du,
\]
where $f_+=\max(f,0)$ and $f_-=\max(-f,0)$.
The pointwise relaxation using the same bounds
$AW\le e^{tu^2}R_m\le E_mW$ has the exact interval for the linear response
\begin{equation}
 [A J_K-(E_m-A)P_-,\ A J_K+(E_m-A)P_+].
 \label{s:pointwise-interval}
\end{equation}
There are explicit numbers $d_-,d_+>0$, given by integrals over
ordered sign intervals, for which
\begin{equation}
 \min C_K\ge-P_-+d_-,\qquad
 \max C_K\le P_+-d_+.
 \label{s:ordered-gains}
\end{equation}
Thus \eqref{s:canonical-interval} strictly improves both endpoints
of the discriminant interval obtained from
\eqref{s:pointwise-interval} with the identical nonlinear payment
$\mathcal B(\overline\rho)$.
\end{theorem}

\begin{proof}
The pointwise relaxation permits an arbitrary measurable multiplier
$0\le a(u)\le1$ in $W\{A+(E_m-A)a\}$.  Choosing $a=1$ exactly
where $f$ is positive or negative proves the two extrema in
\eqref{s:pointwise-interval}; intermediate values follow by convexity.

Since $W>0$, Lemma~\ref{s:kernel-oscillation} supplies a compact
negative interval $I_-$ followed by a compact positive interval $I_+$.
Set
\[
 d_+=\min\left\{\int_{I_-}f_-\,du,
                         \int_{I_+}f_+\,du\right\}>0.
\]
For every $w$, the identity
\[
 P_+-C_K(w)=\int_w^\infty f_+(u)\,du
                              +\int_0^w f_-(u)\,du
\]
is at least $d_+$: before the end of $I_-$ the later positive interval
is wholly omitted, and after that point the negative interval is
wholly included.  This proves the upper inequality.  A positive
interval followed by a negative one gives, by the same argument,
\[
 d_-:=\min\left\{\int_{I_+^{\prime}}f_+\,du,
                         \int_{I_-^{\prime}}f_-\,du\right\}>0,
 \qquad C_K(w)+P_-\ge d_-.
\]
Finally $E_m-A>0$, so both improvements remain strict after adding
the same respective nonlinear errors.
\end{proof}

The right-half-plane condition is material: a reference on the
imaginary axis can have a kernel of one sign.  The strict comparison
just proved can be arbitrarily small and its improved interval may
still contain zero.  It establishes neither a uniform useful gain as
height increases nor optimality among estimators retaining the full
Laplace law.  It compares two specified source relaxations, not all
methods of arithmetic evaluation.

\subsection{Conditional applicability and the remaining source issue}

\begin{proposition}[Local applicability at every putative positive double]
\label{s:tail-local-applicability}
Suppose the unchanged theta field has an ordinary double at $(0,a)$
with $a>0$.  On a sufficiently small fixed proper two-root window about
that event, the unshifted reference \eqref{s:unshifted-head} satisfies
the hypotheses of Theorems~\ref{s:tail-discriminant} and
\ref{s:strict-improvement} for every sufficiently large cutoff $m$.
The assertion is pointwise in the event; it supplies no uniform cutoff
over all heights and no contradiction to the assumed double.
\end{proposition}

\begin{proof}
In the quantitative double-window construction, impose additionally
$R\le a/2$ and then choose the inner radius and time widths with the
same inequalities.  The proof of its positive wall margins is
unchanged.  Every point in the spatial disk has positive real part.
On this fixed compact window the complete theta tails tend uniformly
to zero on both nested walls.  For all sufficiently large $m$, their
absolute bounds are less than half the corresponding wall floors of
$H$.  Rouch\'e's theorem gives two zeros of $F=H-E$ inside each disk,
none in the collar, and an outer reference floor at least half that of
$H$.  In particular the reference roots lie in the right half-plane,
and $|E/F|\le|E|/(|H|-|E|)<1$ on the outer wall.

The explicit sufficient ratio \eqref{s:tail-ratio} also tends to zero.
By \eqref{s:tail-envelope},
\[
 E_m\int_0^\infty W(u)\cosh(Ru)\,du
 \le\frac{E_m}{A}
       \int_0^\infty R_m(u)e^{tu^2}\cosh(Ru)\,du.
\]
The mass ratio is bounded by \eqref{s:tail-mass-ratio}; the complete
tail integral tends to zero uniformly on the compact time interval by
dominated convergence.  The reference floor stays bounded below, so
\eqref{s:tail-ratio} eventually holds.  The cited theorems now apply.
\end{proof}

The reference in \eqref{s:unshifted-head} is essential.  A shifted or
reflection-corrected head changes the residual; the corresponding
connector terms and shifted integration rays must then be restored
before applying a source identity.  The present proof uses the
unshifted head and its complete tail, and makes no claim to resolve
their small signed cancellation.  This distinction is developed in
\cite{tail05}.

There is also a limitation within the weighted-square source class.
Positive weights for which the relevant source sums and integrals
converge leave the same construction with masses $a_n,b_n$ multiplied
by those weights.  Thus positivity of the Laplace
law does not characterize the canonical coefficient sequence.  When
another source agrees through mode $m+1$ and has the same source value
at zero, it has the same $F,W,A,E_m$ and hence the same monotone
response enclosure at any fixed admissible reference window.  Its
actual probability law and response need not agree with theta's.
No assertion that a prescribed window contains a collision of that
other source follows from such a match.  This observation is separate
from the finite envelope-moment control in
Theorem~\ref{s:finite-moment-control}; their preservation hypotheses
are different.

The analytic output is therefore the actual-source enclosure
\eqref{s:canonical-interval}, its strict matched comparison
\eqref{s:ordered-gains}, and its conditional applicability near each
putative ordinary double.  The sign of the full theta discriminant
lower bound remains open.  The proof and control boundaries were
internally checked in \cite{review08}; these statements do not claim
independent specialist validation.
